\lyseis{Β1.7}{Εφεξής και διαδοχικές γωνίες}
\sol{1} $x\mathrm{Ο}z=35^\circ+50^\circ=85^\circ$.
\sol{2} $\mathrm{ΒΟΓ}=130^\circ-70^\circ=60^\circ$.
\sol{3} $360^\circ-(60+75+50+90)^\circ=85^\circ$.
\sol{4} $\omega=180^\circ-(40+65)^\circ=75^\circ$.
\sol{5} $360^\circ:3=120^\circ$.
\sol{6} $24^\circ30'+37^\circ45'=61^\circ75'=62^\circ15'$.
\sol{7} α) Λάθος: χρειάζεται και κοινή πλευρά. β) Σωστό. γ) Σωστό. δ) Σωστό.
\sol{8} $\omega=85^\circ$.
\sol{9} $6$ φέτες: $2\cdot75+3\cdot50=300$. Η έκτη $360-300=60^\circ$.
\sol{10} $60^\circ:12^\circ=5$ φορές και $5\cdot48=240$ λεπτά, δηλαδή $4$ ώρες.
\sol{11} $180^\circ-135^\circ=45^\circ$.
\sol{12} $360^\circ:3=120^\circ$.
\sol{13} Γιατί οι γωνίες γύρω από ένα σημείο καλύπτουν όλον τον πλήρη κύκλο, που είναι $360^\circ$, όσες κι αν είναι.
\sol{14} Ναι. Δύο αμβλείες, π.χ. $100^\circ$ και $100^\circ$, είναι εφεξής και η $x\mathrm{Ο}z=200^\circ$ είναι μη κυρτή.
\sol{15} $x+2x+3x=6x=360^\circ$, $x=60^\circ$, και οι γωνίες είναι $60^\circ$, $120^\circ$, $180^\circ$.
\sol{16} $x\mathrm{Ο}z=60^\circ+20^\circ=80^\circ$. $\delta\mathrm{Ο}y=30^\circ$, άρα $\delta\mathrm{Ο}z=30^\circ+20^\circ=50^\circ$.
\sol{17} Ελεύθερη απάντηση. Το άθροισμα είναι $360^\circ$ (περίπου, λόγω μέτρησης).
\sol{18} Ελεύθερη απάντηση.
